% Part: first-order-logic
% Chapter: completeness
% Section: compactness-direct

\documentclass[../../../include/open-logic-section]{subfiles}

\begin{document}

\iftag{FOL}
      {\olfileid{fol}{com}{cpd}}
      {\olfileid{pl}{com}{cpd}}

\olsection{संहतता प्रमेयाची प्रत्यक्ष सिद्धता}

संपूर्णता प्रमेयाच्या सिद्धतेतील कल्पनाच वापरून, संपूर्णता प्रमेयाचा आधार
न घेता संहतता प्रमेय प्रत्यक्ष सिद्ध करता येते. संपूर्णता प्रमेयाच्या
सिद्धतेत आपण !!{sentence}s च्या सुसंगत संच~$\Gamma$ पासून सुरुवात केली,
त्याचा विस्तार करून !!{sentence}s चा सुसंगत\iftag{FOL}{, संतृप्त}{} आणि
!!{complete} संच~$\Gamma^*$ मिळवला; आणि मग~$\Gamma^*$ पासून रचलेल्या
\iftag{FOL}{पद-प्रतिमान~$\Struct{M(\Gamma^*)}$}{!!{valuation}~$\pAssign{v(\Gamma^*)}$}
मध्ये~$\Gamma$ मधील सर्व !!{sentence}s सत्य आहेत, म्हणजे~$\Gamma$
पूर्ततायोग्य आहे, हे दाखवले.

सांततः पूर्ततायोग्य वाक्यसंच पूर्ततायोग्य असतो हे दाखवण्यासाठी हीच पद्धत
वापरता येते. फक्त सत्यता पूर्वप्रमेयापर्यंत नेणाऱ्या निष्पत्तींच्या अशा
समांतर आवृत्त्या सिद्ध कराव्या लागतील की त्यांत ``सुसंगत'' ऐवजी ``सांततः
पूर्ततायोग्य'' वापरलेले असेल.

\begin{prop}
\ollabel{prop:fsat-ccs}
$\Gamma$ हा !!{complete} आणि सांततः पूर्ततायोग्य आहे असे समजा. मग:
\begin{enumerate}
\tagitem{prvAnd}{$(!A \land !B) \in \Gamma$
  तेव्हा आणि केवळ तेव्हाच, जेव्हा $!A \in \Gamma$ आणि $!B \in \Gamma$
  ही दोन्ही विधाने लागू होतात.}{}

\tagitem{prvOr}{$(!A \lor !B) \in \Gamma$ तेव्हा आणि केवळ तेव्हाच,
  जेव्हा $!A \in \Gamma$ किंवा $!B \in \Gamma$ यांपैकी एक लागू होते.}{}

\tagitem{prvIf}{$(!A \lif !B) \in \Gamma$ तेव्हा आणि केवळ तेव्हाच,
  जेव्हा $!A \notin \Gamma$ किंवा $!B \in \Gamma$ यांपैकी एक लागू होते.}{}
\end{enumerate}
\end{prop}

\tagprob{FOL}
\begin{prob}
\olref[fol][com][cpd]{prop:fsat-ccs} सिद्ध करा. $\Proves$ चा वापर टाळा.
\end{prob}
\tagendprob

\tagprob{notFOL}
\begin{prob}
\olref[pl][com][cpd]{prop:fsat-ccs} सिद्ध करा. $\Proves$ चा वापर टाळा.
\end{prob}
\tagendprob

\iftag{FOL}{%
\begin{lem}
\ollabel{lem:fsat-henkin} प्रत्येक सांततः पूर्ततायोग्य संच~$\Gamma$ चा
विस्तार करून संतृप्त आणि सांततः पूर्ततायोग्य संच~$\Gamma'$ मिळवता येतो.
\end{lem}
}{}

\tagprob{FOL}
\begin{prob}
\olref[fol][com][cpd]{lem:fsat-henkin} सिद्ध करा. (सूचना: महत्त्वाची
पायरी म्हणजे !!{derivation}s किंवा सुसंगततेचा अजिबात आधार न घेता,
$\Gamma_n$ सांततः पूर्ततायोग्य असेल, तर $\Gamma_n \cup \{!D_n\}$
देखील सांततः पूर्ततायोग्य आहे हे दाखवणे.)
\end{prob}
\tagendprob

\iftag{FOL}{
\begin{prop}\ollabel{prop:fsat-instances}
$\Gamma$ हा संपूर्ण, सांततः पूर्ततायोग्य आणि संतृप्त आहे असे समजा.
\begin{tagenumerate}{prvEx,prvAll}
\tagitem{prvEx}{$\lexists[x][!A(x)] \in \Gamma$ तेव्हा आणि केवळ तेव्हाच,
  जेव्हा किमान एका बंद पद~$t$ साठी $!A(t) \in \Gamma$.}{}
\tagitem{prvAll}{$\lforall[x][!A(x)] \in \Gamma$ तेव्हा आणि केवळ तेव्हाच,
  जेव्हा सर्व बंद पदे~$t$ यांच्यासाठी $!A(t) \in \Gamma$.}{}
\end{tagenumerate}
\end{prop}
}{}

\tagprob{FOL}
\begin{prob}
\olref[fol][com][cpd]{prop:fsat-instances} सिद्ध करा.
\end{prob}
\tagendprob

\begin{lem}
\ollabel{lem:fsat-lindenbaum} प्रत्येक सांततः पूर्ततायोग्य संच~$\Gamma$
चा विस्तार करून !!a{complete} आणि सांततः पूर्ततायोग्य संच~$\Gamma^*$
मिळवता येतो.
\end{lem}

\tagprob{FOL}
\begin{prob}
\olref[fol][com][cpd]{lem:fsat-lindenbaum} सिद्ध करा. (सूचना: महत्त्वाची
पायरी म्हणजे $\Gamma_n$ सांततः पूर्ततायोग्य असेल, तर $\Gamma_n \cup
\{!A_n\}$ किंवा $\Gamma_n \cup \{\lnot !A_n\}$ यांपैकी एक संच सांततः
पूर्ततायोग्य असतो हे दाखवणे.)
\end{prob}
\tagendprob

\tagprob{notFOL}
\begin{prob}
\olref[pl][com][cpd]{lem:fsat-lindenbaum} सिद्ध करा. (सूचना: महत्त्वाची
पायरी म्हणजे $\Gamma_n$ सांततः पूर्ततायोग्य असेल, तर $\Gamma_n \cup
\{!A_n\}$ किंवा $\Gamma_n \cup \{\lnot !A_n\}$ यांपैकी एक संच सांततः
पूर्ततायोग्य असतो हे दाखवणे.)
\end{prob}
\tagendprob

\begin{thm}[संहतता]
\ollabel{thm:compactness-direct} $\Gamma$ पूर्ततायोग्य असतो तेव्हा आणि
केवळ तेव्हाच, जेव्हा तो सांततः पूर्ततायोग्य असतो.
\end{thm}

\begin{proof}
$\Gamma$ पूर्ततायोग्य असेल, तर अशी
\iftag{FOL}{!!a{structure}~$\Struct{M}$}{!!a{valuation}~$\pAssign{v}$}
असते की प्रत्येक $!A \in \Gamma$ साठी
$\iftag{FOL}{\Sat{M}}{\pSat{v}}{!A}$. अर्थात, ही
\iftag{FOL}{$\Struct{M}$}{$\pAssign{v}$} देखील~$\Gamma$ चा प्रत्येक
सांत उपसंच पूर्त करते; म्हणून~$\Gamma$ सांततः पूर्ततायोग्य आहे.

आता $\Gamma$ सांततः पूर्ततायोग्य आहे असे समजा.\iftag{FOL}{
  \olref{lem:fsat-henkin} नुसार असा सांततः पूर्ततायोग्य, संतृप्त संच
  $\Gamma' \supseteq \Gamma$ असतो.}{} \olref{lem:fsat-lindenbaum} नुसार
$\iftag{FOL}{\Gamma'}{\Gamma}$ चा विस्तार करून !!a{complete} आणि सांततः
पूर्ततायोग्य संच~$\Gamma^*$ मिळवता येतो\iftag{FOL}{, आणि $\Gamma^*$
  अजूनही संतृप्त असतो}{}. \olref[mod]{defn:termmodel} मधीलप्रमाणे
\iftag{FOL}{पद-प्रतिमान~$\Struct{M(\Gamma^*)}$}{!!{valuation}~$\pAssign{v(\Gamma^*)}$}
रचा.\iftag{FOL}{ लक्षात घ्या की \olref[mod]{prop:quant-termmodel} हे
$\Gamma^*$ सुसंगत आहे (किंवा त्या दृष्टीने !!{complete} वा संतृप्त आहे)
या तथ्यावर अवलंबून नव्हते; ते फक्त $\Struct{M(\Gamma^*)}$ हे बंद पदांनी
आच्छादित आहे या तथ्यावर अवलंबून होते.}{}
सत्यता पूर्वप्रमेयाची सिद्धता (\olref[mod]{lem:truth}) तशीच लागू होते,
जर \olref[ccs]{prop:ccs} च्या संदर्भाऐवजी
\olref{prop:fsat-ccs} चा संदर्भ वापरला\iftag{FOL}{ आणि
\olref[hen]{prop:saturated-instances} च्या संदर्भाऐवजी
\olref{prop:fsat-instances} चा संदर्भ वापरला}{}.
%READERNOTE{OLCOM-006: गोठवलेल्या इंग्रजी स्रोताच्या notFOL शाखेत सत्यता पूर्वप्रमेयासाठी prop:ccs च्या जागी prop:fsat-ccs वापरण्याचा गंतव्य-संदर्भ FOL टॅगच्या आत गेल्यामुळे दिसत नाही. मराठीत समान संदर्भ notFOL शाखेलाही स्पष्टपणे लागू केला असून सर्व गोठवलेले संदर्भ जतन केले आहेत.}
\end{proof}

\tagprob{FOL}
\begin{prob}
\olref[fol][com][cpd]{thm:compactness-direct} च्या सिद्धतेला आवश्यक असलेली
सत्यता पूर्वप्रमेयाची (\olref[fol][com][mod]{lem:truth}) संपूर्ण सिद्धता
लिहा.
\end{prob}
\tagendprob

\tagprob{notFOL}
\begin{prob}
\olref[pl][com][cpd]{thm:compactness-direct} च्या सिद्धतेला आवश्यक असलेली
सत्यता पूर्वप्रमेयाची (\olref[pl][com][mod]{lem:truth}) संपूर्ण सिद्धता
लिहा.
\end{prob}
\tagendprob

\end{document}
